\section*{Chapitre \ref{ch:b2:conjugate-additions} --- Carbonylés conjugués~: Michael et Robinson}

\begin{solution}{exo:b2:conjugate-additions:1}
\begin{itemize}
  \item \ce{CH3MgBr}~: surtout 1,2 (dur, irréversible), 1-méthylcyclohex-2-én-1-ol.
  \item \ce{(CH3)2CuLi}~: 1,4, 3-méthylcyclohexanone.
  \item Thiophénol avec une base~: 1,4 (le thiolate est mou), 3-(phénylsulfanyl)cyclohexanone.
  \item \ce{LiAlH4}~: 1,2, cyclohex-2-én-1-ol.
\end{itemize}
\end{solution}

\begin{solution}{exo:b2:conjugate-additions:2}
Donneurs~: la pentane-2,4-dione, le nitrométhane, le propanedioate de diéthyle (liaison \ce{C-H} acide).
Accepteurs~: le propènenitrile, le propénoate de méthyle, la but-3-én-2-one (\ce{C=C} conjuguée avec
un groupe accepteur).
\end{solution}

\begin{solution}{exo:b2:conjugate-additions:3}
\ce{C4H9Br + 2Li -> C4H9Li + LiBr}, puis \ce{2C4H9Li + CuI -> (C4H9)2CuLi + LiI}, dans l'éther à basse
température sous gaz inerte.
\end{solution}

\begin{solution}{exo:b2:conjugate-additions:4}
Le (3-oxobutyl)propanedioate de diéthyle, \ce{(EtOOC)2CHCH2CH2COCH3}. L'hydrolyse donne l'acide
propanedioïque substitué, qui perd \ce{CO2} par chauffage (\cref{prop:b2:enolates-aldol:decarboxylation})~:
l'acide 5-oxohexanoïque, \ce{HOOCCH2CH2CH2COCH3}.
\end{solution}

\begin{solution}{exo:b2:conjugate-additions:5}
Le diméthylcuprate de lithium dans l'éther à basse température, puis une hydrolyse aqueuse~: \omterm{def:b2:conjugate-additions:addition-modes}{addition
conjuguée}. Le bromure de méthylmagnésium, plus dur, s'additionne surtout en 1,2 et donne
majoritairement le 1-méthylcyclohex-2-én-1-ol.
\end{solution}

\begin{solution}{exo:b2:conjugate-additions:6}
\ce{Et3N} déprotone un peu de \ce{EtSH} pour donner \ce{EtS-}~; le thiolate s'additionne sur le carbone
$\beta$, l'\omterm{def:b2:enolates-aldol:enolate}{énolate} formé prend le proton d'une autre molécule de \ce{EtSH}, ce qui régénère \ce{EtS-}.
Produit~: la 4-(éthylsulfanyl)butan-2-one. Le soufre est gros et polarisable, son doublet haut en
énergie~: c'est un \omterm{def:b2:frontier-orbitals:hard-soft}{nucléophile mou}, sous \omterm{def:b2:frontier-orbitals:control}{contrôle orbitalaire}, qui attaque donc là où la \omterm{def:b2:frontier-orbitals:frontier}{LUMO} est la plus
grande, et l'addition est renversable, ce qui favorise aussi l'adduit 1,4, plus stable.
\end{solution}

\begin{solution}{exo:b2:conjugate-additions:7}
Le cyanure attaque plus vite le carbone du carbonyle, plus positif~: la cyanhydrine est le produit
cinétique. Sa formation est renversable~; à plus haute température, avec le temps, le cyanure est libéré
et s'additionne de nouveau sur le carbone $\beta$, ce qui donne l'adduit qui conserve la liaison $\pi$ de
\ce{C=O}, plus forte~: le produit thermodynamique.
\end{solution}

\begin{solution}{exo:b2:conjugate-additions:8}
Avec la 2-méthylcyclohexane-1,3-dione~: la cétone de Wieland--Miescher (figure de l'\omterm{def:b2:conjugate-additions:robinson}{annélation de
Robinson}). Avec la cyclohexanone~: la 4,4a,5,6,7,8-hexahydronaphtalén-2(3\textit{H})-one, une énone
bicyclique dont la \ce{C=C} aboutit à un carbone de jonction des cycles, sans groupe méthyle angulaire.
\end{solution}

\begin{solution}{exo:b2:conjugate-additions:9}
\ce{CH3COCH2CH2CH2COCH3}~: C2 et C6 sont les carbonyles, cinq carbones de C2 à C6 bornes comprises,
soit une 1,5-dicétone. Couper C3--C4~: donneur, l'\omterm{def:b2:enolates-aldol:enolate}{énolate} de la propanone en C3 (mieux, le 3-oxobutanoate
d'éthyle, dont l'ester est éliminé ensuite par hydrolyse et \omterm{def:b2:enolates-aldol:decarboxylation}{décarboxylation})~; accepteur, la
but-3-én-2-one. Conditions~: 3-oxobutanoate d'éthyle, but-3-én-2-one, éthanolate de sodium en quantité
catalytique dans l'éthanol~; puis acide aqueux et chauffage.
\end{solution}

\begin{solution}{exo:b2:conjugate-additions:10}
L'\omterm{def:b2:enolates-aldol:aldol}{aldolisation} crée une liaison \ce{C-C} mais transforme une liaison $\pi$ de \ce{C=O} en une liaison
\ce{C-O} $\sigma$~; le bilan est faible et deux molécules n'en font qu'une, si bien que $K$ est modeste et
la rétroaldolisation facile. L'\omterm{def:b2:conjugate-additions:michael}{addition de Michael} crée une liaison \ce{C-C} au prix d'une liaison $\pi$
\ce{C=C}, plus faible, et conserve les deux
\ce{C=O}~: elle est nettement favorable. La réaction inverse exigerait que l'\omterm{def:b2:enolates-aldol:enolate}{énolate} de la cétone expulse
l'anion propanedioate stabilisé~; l'équilibre est si déplacé du côté de l'adduit qu'en présence d'une
quantité catalytique de base l'adduit subsiste.
\end{solution}

\begin{solution}{exo:b2:conjugate-additions:11}
\ce{C(COOEt)2(CH2CH2COOCH3)2}. Après la première \omterm{def:b2:conjugate-additions:michael}{addition de Michael}, le carbone central porte encore
un hydrogène entre deux esters, aussi acide qu'auparavant~: la base l'arrache et une seconde addition suit.
\end{solution}

\begin{solution}{exo:b2:conjugate-additions:12}
Par l'\omterm{def:b2:enolates-aldol:enolate}{énolate} le plus substitué (vers le carbone qui porte le méthyle)~: l'octalone dont le groupe méthyle
est sur le carbone de jonction, la 4a-méthyl-4,4a,5,6,7,8-hexahydronaphtalén-2(3\textit{H})-one. Par
l'\omterm{def:b2:enolates-aldol:enolate}{énolate} le moins substitué (du côté du \ce{CH2})~: le groupe méthyle aboutit sur un carbone du cycle voisin
de la jonction. Le premier est favorisé dans des conditions d'équilibration (un alcoolate dans son alcool,
ou la voie de l'\omterm{def:b2:amines:imine}{énamine} avec une \omterm{def:b2:amines:class}{amine secondaire} et un acide), qui donnent l'\omterm{def:b2:enolates-aldol:kinetic-enolate}{énolate thermodynamique}, le
plus substitué (\cref{prop:b2:enolates-aldol:enolate-control}).
\end{solution}

\begin{solution}{pb:b2:conjugate-additions:1}
\textbf{1.} Un cycle cyclohexane portant \ce{C=O} en C1 et en C3 et un groupe méthyle sur C2~; l'hydrogène
porté par C2 est le plus acide.
\textbf{2.} Son départ donne un anion dont la charge est délocalisée sur C2 et sur les deux oxygènes~; pour
la cyclohexanone, un seul oxygène la partage, et l'acidité est trop faible pour être mesurée dans l'eau.
\textbf{3.} Charge sur C2~; \ce{C1=C2} avec \ce{O^-} sur C1~; \ce{C2=C3} avec \ce{O^-} sur C3.
\textbf{4.} Oui~: avec l'hydroxyde, $K = 10^{14.0 - 5.3} = 10^{8.7}$ (l'eau comme acide conjugué, d'après
$K_e$)~; avec l'éthanolate, $10^{15.9-5.3} = 10^{10.6}$. Dans les deux cas, la réaction est totale.
% ledger: ke:25C, pka:ethanol
\textbf{5.} \ce{C1(OH)=C2(CH3)}~: C2 porte encore un hydrogène, tout ce qu'il faut à un \omterm{def:b2:enolates-aldol:tautomerism}{énol}~; le groupe
méthyle ne remplace que le second.
\textbf{6.} Mou~: la charge est répartie sur trois atomes et l'anion est stabilisé. Il attaque le carbone
$\beta$ (\cref{prop:b2:conjugate-additions:regio}).
\textbf{7.} Le C2 de l'\omterm{def:b2:enolates-aldol:enolate}{énolate} se lie à l'extrémité \ce{CH2} de la but-3-én-2-one~; l'\omterm{def:b2:enolates-aldol:enolate}{énolate} formé prend
un proton~: 2-méthyl-2-(3-oxobutyl)cyclohexane-1,3-dione.
\textbf{8.} Chaque nouvel \omterm{def:b2:enolates-aldol:enolate}{énolate} prend un proton à une molécule de dione (ou au solvant) et régénère
ainsi l'\omterm{def:b2:enolates-aldol:enolate}{énolate} donneur~: la base n'est pas consommée.
\textbf{9.} L'\omterm{def:b2:enolates-aldol:enolate}{énolate} est ambident~; avec un électrophile carboné mou, sous \omterm{def:b2:frontier-orbitals:control}{contrôle orbitalaire}, il réagit
par le carbone, et l'adduit \ce{C-C} conserve deux liaisons \ce{C=O}, la combinaison la plus forte~; un
adduit sur \ce{O}, s'il se formait, reviendrait en arrière.
\textbf{10.} Un seul~: C2, lié à C1, à C3, au méthyle et à la chaîne.
\textbf{11.} Le carbone du carbonyle de la chaîne et le carbonyle C1 du cycle (ou C3)~: \ce{C(=O)}--\ce{CH2}--\ce{CH2}--C2--C1,
cinq carbones bornes comprises.
\textbf{12.} Pour consommer toute la dione, le réactif le plus coûteux~; seulement léger, parce que la
but-3-én-2-one est très toxique et polymérise.
\textbf{13.} Les C4 et C6 du cycle (voisins de C3 et de C1), le \ce{CH2} de la chaîne voisin de son carbonyle,
et le \ce{CH3} terminal de la chaîne. C2 n'en porte aucun.
\textbf{14.} L'\omterm{def:b2:enolates-aldol:enolate}{énolate} du \ce{CH3} terminal attaque un carbonyle du cycle~: le cycle fermé compte ce
carbone, le carbone du carbonyle de la chaîne, deux \ce{CH2}, C2 et le carbone du carbonyle attaqué~: six atomes.
\textbf{15.} L'\omterm{def:b2:enolates-aldol:enolate}{énolate} du \ce{CH2} interne de la chaîne fermerait un cycle tendu à quatre atomes. Un \omterm{def:b2:enolates-aldol:enolate}{énolate}
du cycle (C4 ou C6) attaquant le carbonyle de la chaîne ferme aussi un cycle à six atomes, mais ponté~: cet
\omterm{def:b2:enolates-aldol:aldol}{aldol} ne peut pas se déshydrater (la double liaison se trouverait en tête de pont d'un petit système
bicyclique) et revient en arrière, tandis que l'\omterm{def:b2:enolates-aldol:aldol}{aldol} à cycles accolés est drainé par la déshydratation.
\textbf{16.} L'\omterm{def:b2:enolates-aldol:aldol}{aldol} porte \ce{OH} sur l'ancien carbone du carbonyle du cycle, désormais à la jonction~; la
perte d'eau avec un hydrogène de l'ancien carbone du \ce{CH3} (désormais en $\alpha$ de la cétone de la
chaîne) donne la \ce{C=C} conjuguée avec cette cétone.
\textbf{17.} La nouvelle double liaison est conjuguée avec le carbonyle~: une élimination E1cB passant par
l'\omterm{def:b2:enolates-aldol:enolate}{énolate}, favorisée par le chauffage (\cref{prop:b2:enolates-aldol:aldol-equilibrium}).
\textbf{18.} Une cétone saturée et une cétone $\alpha,\beta$-insaturée (une cyclohexénone).
\textbf{19.} Le carbone de jonction qui porte le groupe méthyle~: quatre substituants différents.
\textbf{20.} Non~: la dione est achirale, et ses deux carbonyles sont attaqués également (ils sont images
l'un de l'autre dans un miroir par rapport à la chaîne)~; avec des réactifs achiraux, le produit est
racémique. Des amines chirales comme catalyseurs favorisent l'un des carbonyles et donnent un
énantiomère en excès (volume de troisième année).
\textbf{21.} Dione \ce{C7H10O2}~: \qty{126.155}{g/mol}~; \ce{C4H6O}~: \qty{70.091}{g/mol}~; \ce{C11H14O2}~:
\qty{178.231}{g/mol}.
\textbf{22.} \ce{C7H10O2 + C4H6O -> C11H14O2 + H2O}~: C 11, H 16, O 3 de chaque côté.
\textbf{23.} $10.0/126.155 = \qty{0.07927}{mol} = \qty{79.3}{mmol}$.
\textbf{24.} $0.07927 \times 178.231 \times 0.70 = \boldsymbol{\approx \qty{9.89}{g}}$ de
cétone de Wieland--Miescher.
\end{solution}
